\documentclass[11pt]{article}
\usepackage[a4paper,margin=1in]{geometry}
\usepackage{fontspec}
\usepackage[boldfont=false]{xeCJK}
\setCJKmainfont{FandolSong-Regular.otf}
\usepackage{amsmath}
\usepackage{amssymb}
\usepackage{array}
\usepackage{booktabs}
\usepackage{graphicx}
\usepackage{adjustbox}
\usepackage{enumitem}
\setlist[itemize]{topsep=2pt,partopsep=0pt,parsep=0pt,itemsep=2pt,leftmargin=1.4em}
\setlist[enumerate]{topsep=2pt,partopsep=0pt,parsep=0pt,itemsep=2pt,leftmargin=1.6em}
\usepackage{parskip}
\usepackage{fvextra}
\fvset{breaklines=true,breakanywhere=true,fontsize=\small}
\setlength{\parindent}{0pt}

\begin{document}

\begin{center}
  {\LARGE\bfseries Kinematics}\\[0.35em]
  {\large A Level Physics}
\end{center}
\vspace{0.6em}

\section*{Key definitions}

These five quantities come up in almost every \textbf{kinematics} 运动学 question. Learn the exact words — the examiner gives marks for precise wording.

\begin{itemize}
\item \textbf{distance} 距离 — the total length of the path travelled. A \textbf{scalar} 标量.

\item \textbf{displacement} 位移 — the straight-line distance from the start to the end, with a direction. A \textbf{vector} 矢量.

\item \textbf{speed} 速率 — the rate of change of distance with time. A scalar.

\item \textbf{velocity} 速度 — the rate of change of displacement with time. A vector.

\item \textbf{acceleration} 加速度 — the rate of change of velocity with time. A vector.

\end{itemize}

The unit of speed and velocity is $\text{m s}^{-1}$; the unit of acceleration is $\text{m s}^{-2}$.

A common mistake: \textbf{deceleration} 减速度 just means acceleration in the opposite direction to the velocity. It is not a separate quantity.

\section*{Motion graphs}

Many marks come from reading or drawing motion graphs. Two graphs matter.

\subsection*{Displacement–time graph}

The \textbf{gradient} 斜率 (steepness) of a displacement–time graph at a point gives the velocity at that moment.

\begin{itemize}
\item flat line → the object is at rest.

\item straight sloping line → constant velocity (gradient = velocity).

\item curved line → changing velocity. Draw a \textbf{tangent} 切线 at the point and find its gradient.

\end{itemize}

\par\medskip
\noindent\begin{minipage}{\linewidth}
\centering
\adjustbox{max width=\linewidth, max totalheight=0.42\textheight}{%
  \includegraphics{/Users/joschka/Desktop/Programming/eXercise/output/eXam/handouts/a_level_physics/assets/02-p57-displacement-time-worked-example.pdf}}\\[0.3em]
\emph{Displacement–time graph of a car on a test track}
\end{minipage}\par\medskip

\subsection*{Velocity–time graph}

The gradient of a velocity–time graph gives the acceleration at that moment.

The \textbf{area} between the line and the time axis gives the displacement in that time.

\begin{itemize}
\item flat line → constant velocity (zero acceleration).

\item straight sloping line → \textbf{uniform acceleration} 匀加速 (constant acceleration).

\item curved line → changing acceleration.

\item area above the time axis is positive displacement; area below is negative (the object moved backwards).

\end{itemize}

\par\medskip
\noindent\begin{minipage}{\linewidth}
\centering
\adjustbox{max width=\linewidth, max totalheight=0.42\textheight}{%
  \includegraphics{/Users/joschka/Desktop/Programming/eXercise/output/eXam/handouts/a_level_physics/assets/02-p58-velocity-time-graph.pdf}}\\[0.3em]
\emph{Velocity–time graph — gradient gives acceleration, area gives displacement}
\end{minipage}\par\medskip

\par\medskip
\noindent\begin{minipage}{\linewidth}
\centering
\adjustbox{max width=\linewidth, max totalheight=0.42\textheight}{%
  \includegraphics{/Users/joschka/Desktop/Programming/eXercise/output/eXam/handouts/a_level_physics/assets/02-p58-acceleration-time-graph.pdf}}\\[0.3em]
\emph{Acceleration–time graph derived from the same motion}
\end{minipage}\par\medskip

To find the displacement, split the area into triangles and rectangles, or count grid squares. Area of a triangle is $\tfrac{1}{2} \times \text{base} \times \text{height}$; area of a rectangle is $\text{base} \times \text{height}$.

\section*{The four SUVAT equations}

For motion in a straight line with uniform acceleration, we use five symbols: starting velocity $u$, final velocity $v$, acceleration $a$, displacement $s$, and time $t$. Four equations link them:

\[
v = u + at
\]

\[
s = ut + \tfrac{1}{2} a t^{2}
\]

\[
s = \tfrac{1}{2}(u + v)t
\]

\[
v^{2} = u^{2} + 2as
\]

Each equation uses four of the five symbols. To pick the right one: write down what you know and what you want, then choose the equation with exactly those four.

\subsection*{Where the SUVAT equations come from}

You should be able to get these from the definitions of velocity and acceleration:

\begin{itemize}
\item $v = u + at$ comes from $a = (v - u)/t$, the gradient of the line.

\item $s = \tfrac{1}{2}(u + v) t$ is the area under the line — a \textbf{trapezium} 梯形 with parallel sides $u$ and $v$ and width $t$.

\item $s = ut + \tfrac{1}{2}at^{2}$ comes from putting $v = u + at$ into the area.

\item $v^{2} = u^{2} + 2as$ comes from removing $t$ from the first two.

\end{itemize}

\par\medskip
\noindent\begin{minipage}{\linewidth}
\centering
\adjustbox{max width=\linewidth, max totalheight=0.42\textheight}{%
  \includegraphics{/Users/joschka/Desktop/Programming/eXercise/output/eXam/handouts/a_level_physics/assets/02-p57-suvat-displacement-time-parabola.pdf}}\\[0.3em]
\emph{Displacement–time graph for uniform acceleration — the slope at any point equals the instantaneous velocity}
\end{minipage}\par\medskip

If a question asks "which equation can be found using only the gradient of a velocity–time graph?", the answer is $v = u + at$ (the gradient is the acceleration).

\subsection*{Choosing a positive direction}

Pick a positive direction at the start and keep it. Anything pointing the other way gets a minus sign. For a ball thrown straight up, if "up" is positive: $u$ is positive, $a = -g$ (\textbf{gravity} 重力 pulls down), and at the highest point the displacement is positive but the velocity is zero.

\section*{Free fall under gravity}

When \textbf{air resistance} 空气阻力 can be ignored, an object in \textbf{free fall} 自由落体 has a constant acceleration $g \approx 9.81\ \text{m s}^{-2}$ downwards. This is the same for every mass.

For a ball dropped from rest and falling a distance $h$:

\[
h = \tfrac{1}{2}gt^{2}, \qquad v = gt, \qquad v^{2} = 2gh.
\]

For a ball thrown straight up with speed $u$:

\begin{itemize}
\item greatest height: put $v = 0$ in $v^{2} = u^{2} - 2gh$, giving $h = u^{2}/(2g)$.

\item time to reach the top: put $v = 0$ in $v = u - gt$, giving $t = u/g$.

\item total time to fall back to the start height: $2u/g$ (the motion is \textbf{symmetric} 对称).

\end{itemize}

\subsection*{Experiment to find $g$}

A common method: drop an object from rest, then measure the distance $h$ it falls and the time $t$ it takes. Then

\[
g = \frac{2h}{t^{2}}.
\]

Repeat for several heights and plot $h$ against $t^{2}$. The gradient of the best straight line is $g/2$, so $g$ is twice the gradient. Repeating reduces \textbf{random error} 随机误差. An electronic timer — using \textbf{light gates} 光电门, or a switch the ball hits — removes \textbf{reaction-time} 反应时间 error.

\par\medskip
\noindent\begin{minipage}{\linewidth}
\centering
\adjustbox{max width=\linewidth, max totalheight=0.42\textheight}{%
  \includegraphics{/Users/joschka/Desktop/Programming/eXercise/output/eXam/handouts/a_level_physics/assets/02-p53-free-fall-g-experiment.pdf}}\\[0.3em]
\emph{Experimental set-up for measuring the acceleration due to free fall}
\end{minipage}\par\medskip

\section*{Motion in two directions}

\par\medskip
\noindent\begin{minipage}{\linewidth}
\centering
\adjustbox{max width=\linewidth, max totalheight=0.42\textheight}{%
  \includegraphics{/Users/joschka/Desktop/Programming/eXercise/output/eXam/handouts/a_level_physics/assets/02-p60-projectile-parabola-sprinkler.pdf}}\\[0.3em]
\emph{Water jets from a garden sprinkler trace parabola paths — a real example of projectile motion}
\end{minipage}\par\medskip

When an object moves at constant velocity in one direction (say \textbf{horizontal} 水平) and speeds up in a direction at right angles to it (say \textbf{vertical} 竖直, under gravity), the two motions do not affect each other. Treat each direction on its own, with its own SUVAT equation.

\subsection*{Horizontal throw}

An object thrown horizontally with speed $u_{\text{H}}$ from height $h$, with air resistance ignored:

\begin{itemize}
\item \textbf{horizontal}: constant velocity $u_{\text{H}}$. After time $t$, the horizontal distance is $x = u_{\text{H}} t$.

\item \textbf{vertical}: starts from rest and speeds up downwards at $g$. After time $t$, it has fallen $y = \tfrac{1}{2} g t^{2}$ and has vertical velocity $v_{\text{V}} = g t$.

\end{itemize}

The time to reach the ground depends \textbf{only} on the height $h$, not on $u_{\text{H}}$. Solve $h = \tfrac{1}{2} g t^{2}$ for $t$; then the horizontal \textbf{range} 射程 is $u_{\text{H}} t$.

The horizontal-velocity graph is a flat line at $u_{\text{H}}$. The vertical-velocity graph is a straight line from the origin with gradient $g$.

\subsection*{Projectile at an angle}

A \textbf{projectile} 抛体 thrown at speed $u$ at angle $\theta$ above the horizontal:

\par\medskip
\noindent\begin{minipage}{\linewidth}
\centering
\adjustbox{max width=\linewidth, max totalheight=0.42\textheight}{%
  \includegraphics{/Users/joschka/Desktop/Programming/eXercise/output/eXam/handouts/a_level_physics/assets/02-p60-projectile-at-angle.pdf}}\\[0.3em]
\emph{Projectile launched at angle θ — horizontal and vertical motions are independent}
\end{minipage}\par\medskip

\begin{itemize}
\item horizontal \textbf{component} 分量 of the starting velocity: $u_{\text{H}} = u \cos\theta$ (stays constant during the flight).

\item vertical component of the starting velocity: $u_{\text{V}} = u \sin\theta$ (gets smaller, becomes zero at the top, then grows downwards).

\end{itemize}

At the highest point, $v_{\text{V}} = 0$, but $v_{\text{H}}$ is still $u\cos\theta$. The time to the top is $t_{\text{up}} = u\sin\theta / g$; the total flight time (back to the start height) is $2t_{\text{up}}$.

\par\medskip
\noindent\begin{minipage}{\linewidth}
\centering
\adjustbox{max width=\linewidth, max totalheight=0.42\textheight}{%
  \includegraphics{/Users/joschka/Desktop/Programming/eXercise/output/eXam/handouts/a_level_physics/assets/02-p60-projectile-range-diagram.pdf}}\\[0.3em]
\emph{Range R of a projectile launched from and landing on level ground}
\end{minipage}\par\medskip

\subsection*{Bouncing ball}

When a ball bounces, its velocity–time graph is a set of straight sloping lines (constant $g$) with a sudden jump at each bounce (the velocity flips direction, and gets smaller if some \textbf{energy} 能量 is lost). Add up the times and the distances across the bounces.

\section*{Two objects meeting}

When two objects move along the same line in different ways, write a displacement equation for each. Use the same start time and the same positive direction. Then set the two displacements equal (or set their difference to a given gap).

For a goods train at constant velocity $u_{\text{G}}$ and an express train starting from rest with acceleration $a$, both passing the same point at $t = 0$:

\[
s_{\text{G}} = u_{\text{G}} t, \qquad s_{\text{E}} = \tfrac{1}{2} a t^{2}.
\]

They are level again when $s_{\text{G}} = s_{\text{E}}$, giving $t = 2 u_{\text{G}} / a$.

\section*{Tips for solving problems}

\begin{enumerate}
\item \textbf{Draw a diagram} and mark the positive direction.

\item \textbf{List the SUVAT symbols} with their known and unknown values, including signs.

\item \textbf{Choose the SUVAT equation} with exactly the four symbols you have, plus the one you want.

\item For projectile motion, \textbf{split into horizontal and vertical} SUVAT problems, linked only by the time $t$.

\item Always \textbf{check the units} of your answer, and that its size is sensible.

\end{enumerate}


\end{document}
